%&mex --translate-file=il2-pl





% Zaznaczyć, że "host ip" oraz "host ipv4" to to samo i że to jest
% rozszerzenie i przygotowanie do ipv6.

% zastanowić się nad pomysłem implementacji komendy "iface" łączącej
% interfejsy urządzeń.


\input colordvi
\nopagenumbers

% Czcionki pomocnicze
\font\tfont=plr10 scaled \magstep3
\font\chfont=plb10 scaled \magstep2
\font\subchfont=plb10 scaled \magstep1

% Makro odpowiedzalne za rozdział
\def\chap#1{\bigskip\noindent{\chfont #1}\smallskip}
\def\subchap#1{\bigskip\noindent{\subchfont #1}\smallskip}
\def\propchap#1{\bigskip\noindent{\subchfont #1}\smallskip}
\def\mytt#1;#2.{
\noindent\hfil\hbox to 2.5in{\strut\hfill#2\hfil}\hskip20pt{\tt #1}\hfill\break
}

% Nagłówek
\centerline{\tfont Koszek-Matyja Network Simulator}
\medskip
\centerline{\it Wojciech A. Koszek}
\centerline{\tt wkoszek@FreeBSD.czest.pl}
\smallskip
\centerline{\it Piotr Matyja}
\centerline{\tt peter.matyja@gmail.com}
\bigskip
\bigskip

% -- Treść --

\chap{Wprowadzenie}
\noindent
Głównym powodem do stworzenia programu {\tt KMNSIM} była chęć wiarygodnego
odtworzenia zachowania sieci {\tt Ethernet} w sposób wirtualny, bez
konieczności faktycznej ingerencji w strukturę i konfigurację jakiejkolwiek
fizycznej sieci.

Dzięki programowej emulacji zachowania sieci Ethernet, możliwe
jest symulowanie połączeń między komputerami-hostami, koncentratorami,
przełącznikami oraz routerami.

\smallskip
\Red{2009-06-16: funkcjonalność dot. routerów nie jest jeszcze
zaimplementowana}
\smallskip

\chap{Architektura symulatora}
\noindent
Symulator stworzony został w modułowy sposób. Taka architektura niesie za
sobą wiele zalet, od choćby łatwej możlwości testowania oprogramowania, a na
rozszerzalności kończąc. W początkowej fazie projektowania problemem stała
się sama złożoność możliwej do zasymulowania sieci oraz ilość możliwych do
wykonania połączeń.

Naszym zdaniem niemożliwe byłoby dokładne przetestowanie programu, gdyby
został wykonany on w sposób typowy dla aplikacji z graficznym interfejsem
użytkownika. Zdecydowaliśmy się toteż na skorzystanie z architektury
oprogramowania znanego z systemów rodziny {\tt UNIX}, gdzie główną
funkcjonalność realizuje aplikacja tekstowa, uruchamiania w trybie
wstadowym. Wynik działania programu zapisywany jest do plik (bądź plików),
których poprawność można łatwo zweryfikować. Na podstawie generowanego pliku
interfejs graficzny, uruchamiany jako osobny, niezależny program, jest w
stanie zwizualizować użytkownikowi poszczególne kroki czasowe symulacji.

\smallskip
\Red{Natywny graficzny interfejs użytkownika nie jest jeszcze
zaimplementowany; istnieje jednakże możliwość wykorzystania pakietu
inżynierskiego Graphviz to podejrzenia w postaci graficznej plików DOT,
które reprezentują graf połączeń między elementami sieciowymi}.
\smallskip

\chap{Zasada działania}
\noindent
Sposób w jaki działa program jest dość prosty. Użytkownik tworzy opis sieci
w pliku tekstowym przy pomocy zdefiniowanego języka opisu sieci.  Język ten
zastępuje typowe, graficzne metody konfiguracji. Jest jednocześnie formatem
służącym do zachowywania struktury symulowanej sieci.

Proces konfiguracji--tworzenia pliku wejściowego--rozpoczynamy od
specyfikacji wszystkich elementów sieci. Następnie dokonujemy połączeń
między nimi w kontekście spinania "wirtualnego" kabla. Poprawna konfiguracja
poszczególnych urządzeń oraz posiadania choć jednego aktywnego połączenia
umożliwiają rozpoczęcie procesu symulacji.

Każde urządzenie sieciowe do którego można się podłączyć (host, hub, switch,
router) posiada pewną ilość interfejsów. 

Podstawową jednostką nadawania jest {\tt host}.  Posiada on możliwość
nadawania i odbioru danych z sieci, które skierowane zostały do niego.
Minimalną ilością hostów nadających jest 1.  Host posiada jeden interfejs,
który w celu aktywacji musi zostać skonfigurowany adresem fizycznym MAC,
adresem IP oraz maską sieciową.

Hub jest pośrednikiem w wymianie pakietów między hostami/routerami.  Hub
posiada pewną z góry określoną ilość interfejsów. Owe interfejsy nie mogą
ulec konfiguracji. Są one tylko przekaźnikiem napływających ramek.

\smallskip
\Red{Obecna implementacja zakłada, że zarówno hub jak i switch ma 8 portów; 
Specyfikacja dowolnej ilości portów wymaga kosmetycznych zmian w kodzie
symulatora.}
\smallskip

\noindent 
Rola switcha jest analogiczna do roli huba, z jedyną funkcjonalną różnicą,
że zamiast przekazywać odebrany pakiet na wszystkie porty, switch dokonuje
routingu w warstwie drugiej przekazując pakiet tylko na port, z którego
host/router docelowy będą w stanie odebrać dane.

Połączenie między dwoma urządzeniami możliwe jest po specyfikacji 4
parametrów: nazwy i numeru interfejsu miejsca źródłowego połączenia oraz
nazwy i numeru interfejsu miejsca przeznaczenia.

Na tym etapie proces tworzenia pliku konfiguracyjnego jest ukończony.

Program interpretuje ów plik oraz sprawdza jego poprawność. W przypadku
błędnej składni pliku konfiguracyjnego, program przerywa działanie oraz
informuje użytkownika o popełnionych błędach. W przeciwnym razie program
rozpoczyna symulację. Symulacja polega na wykonaniu wyspecyfikowanej przez
użytkonika ilości kroków, po którym nastąpi koniec symulacji.

Ostateczną fazą działania jest wypisanie wyników symulacji na ekran. Program
wspiera kilka formatów wyjściowych -- wcześniej wspomniany format Dot
pakietu Graphviz, format tekstowy użyteczny do obserwacji faktycznej
{\bf~implementacji} symulatora lub, najbardziej użyteczny, plik tekstowy
obrazujący zachowanie sieci w poszczególnych krokach symulacji.


% Host.
\chap{Opis języka konfiguracyjnego: {\tt host}}
\noindent
Pierwszym elementem potrzebnym do symulacji jest host.  Każdy {\bf host}
posiada kolejkę nadawczo-odbiorczą. W przypadku wysyłania danych umieszcza w
buforze nadawczym dane dot. miejsca źródłowego i przeznaczenia, po czym
oznacza dane jako gotowe do wysłania. Automatycznie wraz ze stworzeniem
hosta {\tt NAZWA} stworzony zostanie interfejs o nazwie {\tt NAZWA} i
numerze {\tt 0} (analogicznie do nazewnictwa interfejsów w większości
systemów operacyjnych).

\smallskip
\Blue{
Określenie miejsca docelowego, a co za tym idzie, późniejsza decyzja o
akceptacji do przetwarzania pakietu dokonywana jest na bazie treści pakietu.
}
\smallskip

\propchap{Proponowana składnia dla dodawania hosta}
\mytt host newhost1 create; Spowoduje dodanie hosta do sieci.
\mytt host mojhost1 remove; Usuwa host z sieci.

\chap{Opis języka konfiguracyjnego: iface}
\noindent
Konfiguracja interfejsu może odbywać się w przypadku, gdy dany interfejs
należy do hosta bądź routera.

\mytt iface NAZWA NUMER mac AA:BB:CC:DD:EE:FF; Konfiguruje ustawienia warstwy fizycznej.
\mytt iface NAZWA NUMER ip IP; Konfiguruje adres IP hosta.
\mytt iface NAZWA NUMER netmask NETMASK; Konfiguruje maskę sieciową hosta.

W przypadku hosta {\tt NUMER} musi być równy 0 -- host posiada tylko jeden
interfejs. Z kolei w przypadku huba/switcha możliwe jest wyspecyfikowanie
innych interfejsów w obrębie dostępnych portów.  Przykładowo, dla hosta {\tt
h1} możliwymi komendami są:

\medskip
\mytt iface h1 0 mac 11:22:33:44:55:66; Konfiguruje ustawienia warstwy fizycznej.
\mytt iface h1 0 ip 192.168.1.1; Konfiguruje adres IP hosta.
\mytt iface h1 0 netmask 255.255.255.0; Konfiguruje maskę sieciową hosta.
\medskip

% Hub.
\chap{Opis języka konfiguracyjnego: hub}
\noindent
Pierwszym elementem styczności hosta z siecią jest {\bf hub}. Hub jest
urządzeniem, które posiada wiele portów Ethernet. Ruch odebrany na jednym z
portów jest przekierowywany na wszystkie inne dostępne porty.  W celu
stworzenia huba w sieci korzystamy z polecenia:
\medskip
\mytt hub NAZWA create; Spowoduje dodanie huba do sieci.
\mytt hub NAZWA remove; Usuwa huba z sieci.
\medskip

\noindent
Na przykład stworzenie hub'a {\tt hu1} możliwe jest dzięki poleceniom:

\medskip
\mytt hub hu1 create; Spowoduje dodanie huba {\tt hu1} do sieci.
\mytt hub hu1 remove; Usuwa huba {\tt hu1} z sieci.
\medskip


% Switch.
\chap{Opis języka konfiguracyjnego: switch}
\noindent
W celu stworzenia switcha w sieci korzystamy z polecenia:

\medskip
\mytt switch NAZWA create; Spowoduje dodanie huba do sieci.
\mytt switch NAZWA remove; Usuwa huba z sieci.
\medskip

\noindent
Stworzenie switchaa {\tt sw1} możliwe jest dzięki poleceniom:

\medskip
\mytt switch sw1 create; Spowoduje dodanie switcha {\tt sw1} do sieci.
\mytt switch sw1 remove; Usuwa switch {\tt sw1} z sieci.
\medskip

\chap{Opis języka konfiguracyjnego: router}
\noindent
Router to host posiadający wiele interfejsów sieciowych o różnych adresach
IP. Router posiada również tablicę routingu czyli spis miejsc osiągalnych
poprzez określony interfejs sieciowy.  Router posiada możliwość
przekierowania danych z jednego interfejsu do drugiego -- dzieje się tak w
przypadku, w którym miejsce przeznaczenia wysłanego z hosta pakietu znajduje
się w tablicy routingu routera.  Oznacza to, że router wie, na który
interfejs przekierować pakiet i właśnie to dokonuje.

\propchap{Proponowana składnia dla dodawania routera}
\mytt router NAZWA create; Stworzenie routera.
\mytt router NAZWA remove; Usunięcie routera.
\mytt router NAZWA route IP NETMASK NUMER-INTERFEJSU; Ustala trasę routingu.

\smallskip
\Red{
Ta część jest jeszcze niezaimplementowana.
}
\smallskip

\noindent
Aktywne elementy sieci, które posiadają możliwość nadawania danych mogą
posiadać przypisane do siebie zadanie. Zadanie służy do wymuszenia pewnego
ustalonego ruchu danych w strukturze sieci i jest zasadniczą właściwością
pracy symulatora. Zadanie zostanie wykonane po rozpoczęciu symulacji.

% Reszta.
\chap{Język konfiguracyjny: pozostałe połączenia}
\noindent
Na koniec pozostała nam najbardziej użyteczna grupa instrukcji, a mianowicie
ta specyfikująca poszczególne parametry dotyczące procesu symulacji. Zmiany
parametrów możemy dokonać poprzez:

\mytt set NAZWA WARTOSC; Ustawia zmienną NAZWA na wartość WARTOSC.

\noindent
Zaimplementowana funkcjonalność to:

\mytt set simtime ILE-KROKOW; Ustawia ilość kroków symulacji (domyślnie:10).

\chap{Notki implementacyjne}
\noindent
Symulator został napisany w języku ANSI C. Wybór na tę technologię padł z
powodu łatwej dostępności nagłówków, implementujących struktury pakietów
sieci Ethernet oraz najlepszą z dostępnych dróg na zapewnienie modularności
(biblioteki dzielone) naszego projektu.

\Blue{
Architektura symulatora w chwili obecnej używa uproszczonych nagłówków
Ethernet, IP oraz ICMP.
Możliwe jest jednak uzyskanie PEŁNEJ zgodności ze strukturą rzeczywistych
pakietów i wykorzystanie replik w symulatorze.
}

\Red{
Faktem jest, że symulator rzeczywiście pracuje w kontekście "wysyłania",
"odbierania", "analizy" i "przetwarzania" rzeczywistych pakietów. Wiadomość
ICMP REQUEST jest faktyczną porcją danych docierającą do interfejsu hosta,
którą ten musi przeanalizować.
}

Plik nagłówkowy {\tt queue.h} został zapożyczony z projektu FreeBSD i służy
do wygodnego zarządzania strukturami danych (listy, kolejki) wykorzystanymi
w projekcie.

{\Red{Do gruntownego rozszerzenia!}}

\chap{Zakończenie}
\noindent Życzymy miłego użytkowania i czekamy na uwagi.
\bigskip
\bigskip
\rightline{Wojciech A. Koszek}
\rightline{\tt wkoszek@FreeBSD.czest.pl}
\smallskip
\rightline{Piotr Matyja}
\rightline{\tt peter.matyja@gmail.com}

\bye
